\documentstyle[%  


righttag%  


,11pt%  


,amssymb%  


,russian%               remove in Eng.  


,izvru%                 Set izven% in Eng.  


]{amsart}  


  


%       Math definitions  


  


\newcommand{\intl}{\int\limits}  


\newcommand{\supp}{\operatorname{supp}}  


\newcommand{\dist}{\operatorname{dist}}  


\newcommand{\conv}{\operatorname{conv}}  


\newcommand{\myfrac}[2]{\frac{\partial #1}{\partial #2}}  


\newcommand{\tts}[1]{}  


  


%\hoffset=-15mm%                  For Eng.  


   


\setcounter{page}{53}  


\god{1998}  


\nomer{10~(437)} %                        Remove in Eng.  


%\nomer{Vol.\,42, No.\,10}%               Set in Eng.  


%\str{\pageref{begin}--\pageref{end}}%  Set in Eng.  


\udk{517.518}  


\author{\it Б.А.\,Кац}  


% NB:   ^^ remove in Eng.  


\title{Интегрирование по плоской фрактальной кривой,\\  


задача о скачке и обобщенные меры}  


  


\thanks{Работа выполнена при поддержке Российского фонда  


фундаментальных исследований (грант \No\,95-01-00674).}  


\date{}  


% \pagestyle{myheadings}%         Set in Eng.  


  


\pagestyle{plain} %              Remove in Eng.  


\begin{document}  


  


% \thispagestyle{plain}%          Set in Eng.  


  


\maketitle  


\label{begin}  


 


К настоящему времени реализовано несколько подходов к построению теории 


интегрирования по фрактальным кривым и по неспрямляемым кривым более 


обширных классов. В данной статье интегрирование 


по фрактальной кривой рассматривается как обобщенная функция или 


как обобщенная мера, т.\,е. мера со значениями в пространстве 


обобщенных функций. При таком подходе задача о фрактальном 


интегрировании оказывается связанной с одной краевой задачей для 


голоморфных функций, а именно с задачей о скачке. 


 


\section{Введение} 


Построение интеграла $\intl_\Gamma f(z)dz$ по фрактальному контуру 


стало в последние годы довольно актуальным вопросом. Так, одним 


из типичных примеров фрактальной кривой является трещина 


(в пластине, плите и т.\,п.; см. \cite{1}, \cite{2}). Известно, что при 


решении краевых задач теории упругости в плоской области с 


прямолинейным разрезом, моделирующей плиту с трещиной, важную роль 


играет интеграл типа Коши по контуру трещины. Поэтому хотелось бы 


уметь вычислять аналогичный интеграл по фрактальной кривой. Но 


такая кривая неспрямляема, и интеграл по ней в классическом смысле не 


существует. Имеется, однако, несколько способов определения 


этого интеграла в некотором обобщенном смысле. 


 


В работах \cite{3}, \cite{4} интеграл определяется следующим 


образом. Пусть $\Gamma$ есть простая замкнутая жорданова кривая 


на комплексной плоскости $\Bbb C$, разбивающая ее на области $D^+$ 


и $D^-$ так, что $\infty\in D^-$, 


и ориентированная положительным образом 


относительно 


$D^+$. Пусть на $\Gamma$ задана непрерывная функция $f(z)$. Если она 


является 


следом на $\Gamma$ непрерывной в $\ov{D^+}$ 


функции $u(z)$ с интегрируемыми первыми производными, то для спрямляемой 


кривой $\Gamma$ формула Стокса дает 


\begin{equation} 


\intl_\Gamma f(z)dz=-\iint\limits_{D^+}\myfrac{u}{\ov z}dz\,d\ov z. 


\end{equation} 


Для неспрямляемой кривой правая часть равенства (1) может служить 


определением левой. Как отмечает Дж.\,Харрисон \cite{5}, это замечание 


восходит к Уитни. При таком определении интеграла его можно 


вычислить, предварительно продолжив плотность $f(t)$ с кривой $\Gamma$ 


в область $D^+$. Здесь возникает вопрос о существовании продолжения с 


требуемыми свойствами и о возможной зависимости интеграла 


от этого продолжения. В работах \cite{3}, \cite{4} этот вопрос решается 


в терминах условия Г\"ельдера. При $A\subset\Bbb C$, $0<\nu\le 1$ класс 


Г\"ельдера $H_\nu(A)$ по определению состоит из непрерывных на 


множестве $A$ функций $f(t)$, для которых величина 


\begin{equation} 


 h_\nu(f,A)\equiv\sup\bigg\{ 


\frac{|f(t_1)-f(t_2)|}{|t_1-t_2|^\nu}:t_{1,2}\in A,\ t_1\ne t_2 


\bigg\} 


\end{equation} 


конечна. В \cite{3}, \cite{4} показано, что при условии 


\begin{equation} 


\nu>d-1,


\end{equation} 


где $d$ --- клеточная размерность множества $\Gamma$ (она же --- его 


верхняя 


метрическая размерность, box dimension 


и т.\,п., см. \cite{1}, \cite{2}), продолжение Уитни 


функции $f\in H_\nu(\Gamma)$ имеет интегрируемые в $\ov{D^+}$ 


производные и 


правая часть (1) не зависит от выбора конкретного продолжения Уитни 


$u(z)$ (о конструкции продолжений Уитни см., напр., в \cite{6}, 


гл.\,4). 


 


Для определения интеграла по разомкнутой дуге $\Gamma$ в \cite{4} 


предложен следующий прием. Пусть $a$ и $b$ --- начало и конец 


дуги $\Gamma$. Если эти точки можно соединить спрямляемой дугой 


$\gamma$, не 


имеющей с $\Gamma$ иных общих точек (в этом случае дуга $\Gamma$ названа 


в \cite{4} 


достижимой), то $\Gamma\cup\gamma$ 


есть простая замкнутая кривая, интегрирование по 


которой определено формулой (1), и можно положить 


\begin{equation} 


\intl_\Gamma f(z)dz=\intl_{\Gamma\cup\gamma}f(z)dz- 


\intl_\gamma f(z)dz.


\end{equation} 


В \cite{4}, \cite{5}, \cite{7} формула (1) исследуется с позиций  


геометрической теории интеграла. Если функция $f$ первоначально задана 


во всей плоскости $\Bbb C$, то величину интеграла 


$\intl_\Gamma f(z)dz$ можно 


рассматривать как функцию, заданную на некотором множестве кривых, 


например, на множестве ломаных. На этом множестве задается метрика. 


Если интеграл непрерывен в этой метрике, то он продолжим по  


непрерывности на замыкание множества ломаных. В \cite{4}, \cite{5} 


\cite{7} предполагается, что $f\in H_\nu(\Gamma)$, 


а метрика на множестве ломаных 


строится таким образом, чтобы соответствующее замыкание содержало все 


$d$-мерные кривые при заданном $d$, $1<d\le 2$. Выясняется, что интеграл 


непрерывен в такой метрике при условии (3). 


Наиболее общее описание этих результатов дается в \cite{5}. Отметим, 


что вышеописанное продолжение интеграла реализуется все той 


же формулой (1). 


 


В \cite{8} рассматривается определение интеграла по плоской 


разомкнутой неспрямляемой дуге с помощью конструкции, двойственной 


по отношению к только что описанному геометрическому продолжению. 


Значения интеграла от голоморфной функции $f$ по дуге с началом $a$ 


и концом $b$ определяются здесь по формуле Ньютона--Лейбница 


$\intl_\Gamma f(z)dz=F(b)-F(a)$, $F'(z)=f(z)$. 


Тем самым на линеале голоморфных 


вблизи $\Gamma$ функций задается функционал 


$\intl_\Gamma f(z)dz$, выясняются условия, при 


которых этот функционал непрерывен в норме 


пространства Г\"ельдера $H_\nu(\Gamma)$ и, следовательно, может быть 


продолжен на 


замыкание вышеупомянутого линеала в $H_\nu(\Gamma)$. В \cite{8} показано, 


что 


искомое продолжение интеграла имеет вид 


\begin{equation} 


\intl_\Gamma f(z)dz=-\iint\limits_{\Bbb C} 


\myfrac{u}{\ov z}k_\Gamma (z)dz\,d\ov z, 


\end{equation} 


где $u(z)$ есть продолжение $f$ с $\Gamma$ в $\Bbb C$, имеющее 


интегрируемые частные производные и 


компактный носитель, а $k_\Gamma(z)$ --- однозначная ветвь функции 


\begin{equation} 


k_\Gamma (z)=\frac{1}{2\pi i}\ln\frac{z-a}{z-b}, 


\end{equation} 


выделенная посредством разреза по $\Gamma$ и условия 


$k_\Gamma(\infty)=0$. Формула (5) 


установлена в \cite{8} в предположении, что ядро 


$k_\Gamma$ интегрируемо вблизи 


дуги $\Gamma$ в достаточно высокой степени. В \cite{9} 


определение интеграла (1) используется в ситуации, когда плотность 


$f$ не 


является, вообще говоря, непрерывной функцией, но 


принадлежит пространству О.В.\,Бесова. Ее продолжение $u(z)$ 


получается с помощью оператора продолжения типа Уитни, построенного 


в \cite{10}. 


 


Наконец, Т.Е.\,Объедков \cite{11} дал определение интеграла по 


неспрямляемой кривой средствами нестандартного анализа. Его  


стандартизация также приводит к формуле (1). Таким образом, все 


известные 


к настоящему времени версии обобщенного интегрирования по неспрямляемой 


кривой сводятся к формулам (1), (4), (5). Отметим, что формула (5) в 


отличие от (4) справедлива и в ситуации, когда дуга $\Gamma$ 


не является 


достижимой. Если же она достижима, то применение этих двух формул 


приводит к одинаковому результату. 


 


Функционал $\intl_\Gamma f(z)dz$ оказывается при этом ограниченным по 


норме пространств Г\"ельдера или Бесова, но не пространства непрерывных 


функций. Вероятно, ограниченность функционала (1) или (4) в 


норме пространства $C(\Gamma)$ влечет спрямляемость кривой $\Gamma$. 


Прямое 


доказательство этого предположения было получено Т.Е.\,Объедковым в 


его неопубликованной студенческой работе для случая, когда $\Gamma$ 


представляет собою график вещественной функции. 


 


\section{Интегрирование по фрактальной кривой как обобщенная функция} 


 


Как обычно, через $C^\infty(D)$ обозначаем множество всех заданных в 


области $D\subset\Bbb C$ 


функций, имеющих частные производные любого порядка, а через 


$C^\infty_0(D)$ --- множество тех из них, которые имеют 


компактные носители в $D$. При $D=\Bbb C$ эти множества обозначаются 


через $C^\infty$ и $C^\infty_0$ соответственно. 


 


Согласно классической теории распределений любое распределение $\varphi$ 


с носителем на $\Gamma$ может рассматриваться как 


интегрирование по $\Gamma$ некоторой обобщенной функции $f$


$$ 


\varphi(\omega)=\int_\Gamma f(z)\omega(z)dz,\quad 


\omega\in C^\infty. 


$$ 


Если, в частности, $\omega(z)=1$ 


на $\Gamma$, то $\varphi(\omega)$ можно рассматривать как интеграл от 


функции по $\Gamma$.


 


Поэтому интегрирование по фрактальной кривой $\Gamma$ 


можно определить как 


обобщенную функцию с носителем $\Gamma$, обладающую какими-либо 


дополнительными свойствами, позволяющими восстановить соответствующую 


плотность $f$. Чтобы охарактеризовать эти свойства, рассмотрим 


преобразования Коши распределений с носителями на 


$\Gamma$. Преобразование 


Коши распределения $\varphi$, $\supp\varphi\subset\Gamma$, 


есть определенная в $\Bbb C\setminus\Gamma$ функция 


$$ 


k_\varphi(z)=(2\pi i)^{-1}\varphi((\cdot-z)^{-1}); 


$$ 


разумеется, здесь распределение $\varphi$ применяется к некоторой 


функции $\omega(t)\in C^\infty_0$, совпадающей с $(t-z)^{-1}$ вблизи 


$\Gamma$. Очевидно, 


$k_\varphi(z)$ есть голоморфная в $\ov{\Bbb C}\setminus\Gamma$ 


функция и $k_\varphi(\infty)=0$. Как 


обычно, мы отождествляем $k_\varphi$ 


с распределением 


$$


k_\varphi(\omega)=\iint\limits_{\Bbb C} k_\varphi 


(x+iy)\omega(x+iy)dx\,dy,\quad 


\omega\in C^\infty_0(\Bbb C\setminus\Gamma). 


$$


 


Далее, рассмотрим распределение $E=(\pi z)^{-1}$ и свертку 


$\psi=\varphi* E$; 


она существует, поскольку носитель $\varphi$ компактен. Пусть носитель 


функции $\omega\in C^\infty_0$ лежит в $\Bbb C\setminus\Gamma$. 


Несложные преобразования дают (\cite{12}, гл.\,4) 


$$ 


\psi(\omega)=-2i\iint\limits_{\Bbb C}k_\varphi(x+iy) 


\omega(x+iy)dx\,dy=-2i k_\varphi(\omega). 


$$ 


Таким образом, $-2i k_\varphi$ есть сужение $\psi$ на 


$\Bbb C\setminus\Gamma$. 


Если $\Gamma$ имеет 


нулевую площадь и функция $k_\varphi$ интегрируема вблизи $\Gamma$, то 


распределение 


$-2i k_\varphi$ продолжимо на всю плоскость $\Bbb C$, и на 


$C^\infty_0$ определено 


распределение 


\begin{equation} 


s_\varphi=\varphi*E+2i k_\varphi. 


\end{equation} 


Очевидно, $\supp s_\varphi\subset\Gamma$. 


Вообще говоря, $s_\varphi\ne 0$; соответствующий 


пример дает распределение $\varphi=\myfrac{\delta_a}{\ov z}$, 


$a\in\Gamma$. 


 


\begin{defn} Будем называть распределение $\varphi$ с носителем 


на $\Gamma$ регулярным, если преобразование Коши этого распределения 


интегрируемо в любой конечной части плоскости и $s_\varphi=0$. 


\end{defn} 


 


Для регулярного распределения $\varphi$ имеем $\varphi*E=-2i k_\varphi$. 


Но $\myfrac{E}{\ov z}=\delta_0$, т.\,е.{} 


$\myfrac{}{\ov z}(\varphi*E)=\varphi*\delta_0=\varphi$, 


и $\varphi=-2i\myfrac{k_\varphi}{\ov z}$. 


Итак, регулярное распределение действует по формуле 


\begin{equation} 


\varphi(\omega)=2i\iint\limits_{\Bbb C} \myfrac{\omega}{\ov z} 


k_\varphi(x+iy)dx\,dy=-\iint\limits_{\Bbb C}\myfrac{\omega}{\ov z} 


k_\varphi(z)dz\,d\ov z. 


\end{equation} 


 


Если $\Gamma$ --- спрямляемая кривая и $f\in L^1(\Gamma,ds)$, 


то интегрирование $\varphi(\omega)=\intl_\Gamma f(z)\omega(z)dz$ 


является регулярным распределением. Его преобразование Коши есть 


интеграл типа Коши 


$$ 


k_\varphi(z)=\frac{1}{2\pi i}\int_\Gamma\frac{f(t)dt}{t-z}. 


$$ 


В почти всех (относительно $ds$) точках $t\in\Gamma$ эта функция имеет 


некасательные предельные значения с обеих сторон, связанные 


соотношением 


\begin{equation} 


k^+_\varphi(t)-k^-_\varphi(t)=f(t). 


\end{equation} 


Если $f\in H_\nu(\Gamma)$ и $\nu>1/2$, 


то функция $k_\varphi$ имеет непрерывные 


граничные значения на всей кривой $\Gamma$ 


за возможным исключением концов $\Gamma$, 


если это разомкнутая дуга (\cite{13}). Если кривая 


кусочно-гладкая, то это справедливо при любом $\nu$ (напр., \cite{14}, 


\cite{15}). Разумеется, в обоих этих случаях соотношение (9)  


выполняется во всех точках кривой $\Gamma$ за возможным исключением ее 


концов. 


Это наблюдение приводит к следующему определению интегрирования. 


 


\begin{defn} Будем называть регулярное распределение $\Gamma$ 


с носителем на $\Gamma$ интегрированием, если его преобразование Коши 


имеет предельные значения с обеих сторон в каждой точке $\Gamma$ (за 


возможным исключением концов $\Gamma$, если это разомкнутая дуга). 


Разность (9) этих предельных значений есть плотность данного  


интегрирования.


\end{defn} 


 


Обсудим вопрос о существовании интегрирований на неспрямляемых кривых. 


Прежде всего отметим, что формула (1) определяет интегрирование по 


замкнутой неспрямляемой кривой с единичной плотностью. Его 


преобразование Коши есть характеристическая функция $\chi_\Gamma(z)$ 


области 


$D^+$, равная единице в этой области и нулю вне ее. Аналогично, 


формула (5) определяет интегрирование с единичной плотностью по 


разомкнутой дуге, а формула (6) --- его преобразование Коши. 


 


В общем случае вопрос о существовании интегрирования с заданной 


плотностью есть вопрос о разрешимости краевой задачи о восстановлении 


голоморфной в $\ov{\Bbb C}\setminus\Gamma$ 


функции по заданному скачку на кривой 


$\Gamma$, т.\,е. по 


краевому условию (9). Эта задача исследовалась 


автором в работах \cite{16}, \cite{17}. 


 


Там установлено, что градиент продолжения Уитни $f^w(z)$ 


функции $f\in H_\nu(\Gamma)$ 


(в обозначениях \cite{6} $f^w\equiv\cal E_0 f$) интегрируем вблизи 


$\Gamma$ в любой степени, меньшей величины 


$p_1(\nu,d)=\frac{2-d}{1-\nu}$. При 


\begin{equation} 


%%% k(z)=u(z)-\frac{1}{2\pi i} \iint\limits_{\Bbb C} 


%%% \frac{d\zeta\,d\ov\zeta}{\zeta-z} 


\nu>d/2


\end{equation} 


эта величина больше двух, и функция 


\begin{equation} 


k(z)=u(z)-\frac{1}{2\pi i} \intl_{\Bbb C}


\myfrac{u}{\ov\zeta}\frac{d\zeta\,d\ov\zeta}{\zeta-z}


\end{equation} 


при $u(z)=f^w(z)\chi_\Gamma(z)$ 


дает решение краевой задачи (9) на замкнутой кривой 


$\Gamma$. Полагая в (8) $k_\varphi=k$, получаем интегрирование с 


плотностью $f$. На разомкнутой дуге $\Gamma$


формула (11) также дает решение задачи о скачке (9), если положить 


в ней $u(z)=f^w(z)k_\Gamma(z)$, где 


$k_\Gamma$ --- логарифмическое ядро (5), и


$f^w$ должно иметь компактный носитель, т.\,е. $f^w=\xi\cal E_0f$, а 


функция $\xi\in C^\infty_0$ равна единице вблизи $\Gamma$. Здесь, однако, 


возникают 


дополнительные условия разрешимости краевой задачи (9), связанные с 


тем, что на показатель интегрируемости 


$\myfrac{u}{\ov\zeta}=k_\Gamma(\zeta)\myfrac{f^w}{\ov\zeta}$ 


влияет логарифмическое ядро $k_\Gamma$. 


Если $k_\Gamma$ интегрируемо вблизи точек $a$ 


и $b$ в степени $p$, то интеграл (11) имеет смысл при 


\begin{equation} 


\nu>1-(1-p^{-1})(2-d),


\end{equation} 


а определяемая им функция $k(z)$ имеет непрерывные предельные 


значения с обеих сторон на $\Gamma\setminus\{a,b\}$ при условии (10). 


Если концы фрактальной дуги $\Gamma$ можно соединить гладкой дугой 


$\gamma$, не 


имеющей с $\Gamma$ общих внутренних точек, то $k_\Gamma$ 


имеет вблизи точек $a$ 


и $b$ логарифмическую асимптотику. В таком случае показатель 


$p$ в условии (12) можно устремить к бесконечности, и это условие 


совпадает с (3). 


 


Таким образом, справедлива 


 


\begin{thm} Пусть $\Gamma$ есть замкнутая кривая клеточной 


размерности $d$, $f\in H_\nu(\Gamma)$ 


и выполнено условие $(10)$. Тогда существует интегрирование по кривой 


$\Gamma$ с плотностью $f$. На разомкнутой дуге $\Gamma$ 


размерности $d$ интегрирование с заданной плотностью 


$f\in H_\nu(\Gamma)$ 


существует, если кроме $(10)$ выполнено условие $(12)$. 


\end{thm} 


 


\begin{note} Недавно Т.Е.\,Объедков доказал средствами 


нестандартного 


анализа, что на самоподобном фрактале функция (11) имеет непрерывные 


граничные значения на $\Gamma$ с обеих сторон не только при условии 


(10), 


но и при более слабом условии (3). Это значит, что на 


замкнутом самоподобном фрактале $\Gamma$ размерности $d$ интегрирование 


с заданной плотностью $f\in H_\nu(\Gamma)$ существует при условии (3), 


а на разомкнутой --- при условии (12). 


\end{note} 


 


Теперь рассмотрим вопрос о продолжениях интегрирований и об их 


порядках сингулярности. Обычно \cite{12} порядком сингулярности 


распределения $\varphi$ 


называют такое число $n$, что 


$$ 


|\varphi(\omega)|\le B\|\omega\|_{C^n(A)}\ \; 


\forall\omega\in C^\infty, 


$$ 


где положительная постоянная $B$ не зависит от $\omega$, а множество 


$A$ содержит $\supp\varphi$. Согласно (8) порядок сингулярности любого 


регулярного 


распределения не превосходит единицы, и поэтому любое такое 


распределение продолжимо до непрерывного функционала на $C^1$. 


Как 


мы отмечали выше, интегрирование по неспрямляемой кривой не может, 


по-видимому, быть распределением нулевого порядка (т.\,е.  


распределением типа меры). Однако покажем сейчас, что порядок 


сингулярности такого интегрирования все же ниже единицы. 


 


Чтобы говорить о дробных порядках сингулярности, мы должны 


рассмотреть какую-нибудь шкалу пространств $A_\lambda$, 


$0\le\lambda\le 1$, таких, что при $\lambda\to 0$ 


пространство $A_\lambda$ приближается по своим свойствам к 


$C(\Gamma)$, а при $\lambda\to 1$ --- 


к $C^1(\Gamma)$. В такой ситуации порядком сингулярности распределения 


$\varphi$ 


относительно шкалы $A_\lambda$ можно назвать 


точную нижнюю грань таких $\lambda$, что 


$$ 


|\varphi(\omega)|\le B_\lambda\|\omega\|_{A_\lambda}\ \; 


\forall\omega\in C^\infty_0, 


$$ 


где $B_\lambda>0$ не зависит от $\omega$. 


 


Найдем порядок сингулярности интегрирования (8) относительно 


шкалы Г\"ельдера $H_\lambda$. Напомним, что норма в пространстве 


Г\"ельдера 


$H_\lambda(\Gamma)$ определяется равенством 


$$ 


\|\omega\|_{H_\lambda}=h_\lambda(\omega,\Gamma)+ 


\sup\{|\omega(z)|:z\in\Gamma\} 


$$ 


(см. (2)). 


 


\begin{lem} Пусть $\Gamma$ --- 


замкнутая кривая размерности $d<2$, $\omega\in C^\infty_0$, 


а $\wt\omega$ --- продолжение 


Уитни сужения $\omega|_\Gamma$ в $\Bbb C$ с компактным носителем, т.\,е. 


$\wt\omega(z)=\xi(z)\cal E_0(\omega|_\Gamma)(z)$ 


в обозначениях \cite{6}, где $\xi$ есть равная $1$ вблизи $\Gamma$ 


функция класса $C^\infty_0$. Если голоморфная в 


$\Bbb C\setminus\Gamma$ функция $K(z)$ 


ограничена в $\Bbb C$, то 


\begin{equation} 


\iint\limits_{\Bbb C}\myfrac{\omega}{\ov z}K(z)dz\,d\ov z= 


\iint\limits_{\Bbb C}\myfrac{\wt\omega}{\ov z}K(z)dz\,d\ov z. 


\end{equation} 


\end{lem} 


 


\begin{pf} 


Продолжение Уитни 


сохраняет модуль непрерывности, 


так что $\wt\omega\in H_1(\Bbb C)$. 


Но тогда разность $w(z)=\omega(z)-\wt\omega(z)$ также принадлежит 


$H_1(\Bbb C)$, имеет компактный носитель, ограниченные 


первые производные и исчезает на $\Gamma$. Разобьем плоскость на 


квадраты со стороной $\varepsilon>0$ и обозначим через 


$\Delta_\varepsilon$ объединение 


тех из них, которые пересекают $\Gamma$. 


При $d<2$ площадь $\Delta_\varepsilon$ стремится к нулю 


при $\varepsilon\to 0$, 


и поэтому 


$$


\iint\limits_{\Bbb C}\myfrac{w}{\ov z}K(z)dz\,d\ov z= 


\lim_{\varepsilon\to 0} 


\iint\limits_{\Bbb C\setminus\Delta_{\scriptstyle\varepsilon}} 


\myfrac{w}{\ov z}K(z)dz\,d\ov z= 


\lim_{\varepsilon\to 0}\int\limits_{\partial\Delta_{\scriptstyle\varepsilon}}


w(z)K(z)dz=\lim_{\varepsilon\to 0}


\sum_{a\subset\Delta_{\scriptstyle\varepsilon}}


\intl_{\partial Q}w(z)K(z)dz, 


$$


где $Q$ 


--- входящие 


в $\Delta_\varepsilon$ квадраты рассматриваемого разбиения, а 


$\partial\Delta_\varepsilon$ 


и $\partial Q$ --- обычным образом ориентированные границы 


$\Delta_\varepsilon$ и 


$Q$ соответственно. Каждый квадрат $Q$ содержит точку $z_Q\in\Gamma$ и 


$w(z_Q)=0$. Поскольку $|z-z_Q|\le \sqrt{2}\,\varepsilon$ 


при $z\in Q$, то 


$$ 


\bigg|\intl_{\partial Q}w(z)K(z)dz\bigg| 


=\bigg|\intl_{\partial Q}(w(z)-w(z_Q))K(z)dz\bigg| 


\le 4\sqrt{2}\,\varepsilon^2B_K h_1(\omega,\Gamma), 


$$ 


где $B_K$ --- верхняя граница $K$. Пусть $m_\varepsilon$ есть число 


входящих в $\Delta_\varepsilon$ 


квадратов, тогда 


$$ 


\bigg|\sum_{Q\subset\Delta_{\scriptstyle\varepsilon}}\intl_{\partial Q} 


w(z)K(z)dz\bigg|\le B m_\varepsilon\varepsilon^2,\quad 


B=4\sqrt{2}\,h_1(w,\Gamma)B_K. 


$$ 


Но $m_\varepsilon\varepsilon^2$ есть площадь $\Delta_\varepsilon$, 


и она стремится к нулю при $\varepsilon\to 0$. 


\end{pf} 


 


\begin{thm} Пусть $\Gamma$ есть замкнутая кривая размерности $d$, 


а $\varphi$ --- заданное на $\Gamma$ интегрирование. Тогда порядок 


сингулярности $\varphi$ 


в шкале Г\"ельдера не превосходит $d-1$. 


\end{thm} 


 


\begin{pf} Пусть $\omega\in C^\infty_0$, $0<\lambda\le 1$. Очевидно,


$\omega|_\Gamma\in H_\lambda(\Gamma)$. 


Тогда продолжение Уитни $\cal E_0(\omega|_\Gamma)$ 


принадлежит пространству $H_\lambda(\Bbb C)$, 


причем 


$h_\lambda(\cal E_0(\omega|_\Gamma),\Bbb C)=h_\lambda(\omega,\Gamma)$ 


\cite{6}. 


Следовательно \cite{6},


$$ 


\bigg|\myfrac{}{\ov z}\cal E_0(\omega|_\Gamma)\bigg|


\le h_\lambda(\omega,\Gamma)(\dist(z,\Gamma))^{\lambda-1}\le


\|\omega\|_{H_\lambda}(\dist(z,\Gamma))^{\lambda-1}. 


$$ 


При построении функции $\wt\omega(z)=\xi(z)\cal E_0(\omega|_\Gamma)(z)$ 


сомножитель $\xi\in C^\infty_0$ 


можно выбрать так, чтобы любая точка носителя $\xi$ 


была удалена от $\Gamma$ не более, чем на 1. Если при этом 


$\Big|\myfrac{\xi}{\ov z}\Big|\le B$, $B\ge 1$,


то 


$$ 


\bigg|\myfrac{\wt\omega}{\ov z}\bigg|=\bigg| 


\myfrac{\xi}{\ov z}\cal E_0(\omega|_\Gamma)+\xi 


\myfrac{\cal E_0(\omega|_\Gamma)}{\ov z}\bigg| 


\le B\|\omega\|_{H_\lambda}(\dist(z,\Gamma))^{\lambda-1}. 


$$ 


Согласно лемме 1 теперь получим 


\begin{equation} 


|\varphi(\omega)|\le B\|\omega\|_{H_\lambda} 


\iint\limits_{\supp\xi}|k_\varphi(z)|(\dist(z,\Gamma))^{\lambda-1} 


|dz\,d\ov z|, 


\end{equation} 


где $B$ зависит от $\xi$, но не от $\omega$. Поскольку функция 


$k_\varphi$ 


здесь ограничена, то доказательство сводится к вопросу: при каких 


значениях $p>0$ функция $(\dist(z,\Gamma))^{-p}$ интегрируема вблизи 


$\Gamma$? В работе \cite{16} показано, что это так при $p<2-d$.


Неравенство $1-\lambda<2-d$ равносильно условию $\lambda>d-1$. 


\end{pf} 


 


\begin{cor} Пусть $\Gamma$ есть замкнутая кривая размерности $d$. 


Любое заданное на $\Gamma$ 


интегрирование $\varphi$ продолжимо по непрерывности на пространства 


Г\"ельдера 


$H_\lambda(\Gamma)$ при $\lambda>d-1$. Продолженный 


функционал определяется формулой 


\begin{equation} 


\varphi(g)=-\iint\limits_{\Bbb C}\myfrac{g^w}{\ov z}k_\varphi(z)


dz\,d\ov z,\quad g\in H_\lambda(\Gamma), 


\end{equation} 


где $g^w=\xi\cal E_0g$, $\cal E_0$ --- оператор продолжения Уитни 


с множества $\Gamma$ в плоскость $\Bbb C$, а $\xi\in C^\infty_0$ и 


$\xi(z)=1$ вблизи $\Gamma$. 


\end{cor} 


 


Доказательство немедленно вытекает из теоремы 2 и непрерывности 


оператора Уитни $\cal E_0$.


 


Таким образом, хотя интегрирование существует, вообще говоря, 


лишь для плотности $f$, удовлетворяющей условию Г\"ельдера с показателем 


$\nu>d/2$, оно позволяет вычислить интеграл 


$\varphi(g)$\; $\Big(=\intl_\Gamma f(t)g(t)dt\Big)$ 


при любом $g\in H_\lambda(\Gamma)$, $\lambda>d-1$. Если $f(t)\ne 0$ на 


$\Gamma$, то 


произведение $fg$ пробегает все пространство $H_\lambda(\Gamma)$. 


 


Если $\Gamma$ --- разомкнутая дуга, то преобразование Коши 


интегрирования 


по $\Gamma$ уже не обязано быть ограниченным. На концах дуги 


оно может иметь особенности. Однако доказательство леммы 1 нетрудно 


перенести на случай, когда функция $K(z)$ ограничена вне произвольной 


окрестности концов дуги $\Gamma$ и интегрируема внутри этой 


окрестности. Оценивая правую часть (14) с помощью неравенства 


Г\"ельдера, нетрудно убедиться, что функционал $\varphi$ ограничен по 


норме 


$H_\lambda$ 


при условии 


\begin{equation} 


\lambda>1-(2-d)(1-p^{-1})=d-1+(2-d)p^{-1}, 


\end{equation} 


где $p$ --- показатель интегрируемости преобразования Коши $k_\varphi$. 


В частности, если функция $k_\varphi$ интегрируема вблизи концов 


$\Gamma$ в 


любой 


степени, то получим ту же оценку для порядка сингулярности 


$\varphi$, что и в теореме 2. 


 


Теперь перейдем к вопросам единственности интегрирований и 


интегралов. Прежде всего решим вопрос о единственности интегрирования с 


заданной плотностью, или, иначе говоря, о существовании 


нетривиальных интегрирований с нулевой плотностью. 


 


Преобразование Коши интегрирования по замкнутой кривой $\Gamma$ с 


нулевой плотностью представляет собой непрерывную в $\ov{\Bbb C}$ и 


голоморфную в 


$\ov{\Bbb C}\setminus\Gamma$ функцию, обращающуюся в нуль в точке 


$\infty$. Если $\Gamma$ --- спрямляемая кривая, то такая функция есть 


тождественный нуль 


по теореме Пенлеве (см., напр., \cite{18}, гл.\,3). Это означает, 


что интегрирование по спрямляемой кривой единственным образом 


определяется 


своей плотностью, т.\,е. на спрямляемой кривой нет никаких  


интегрирований, кроме классических. Если же 


$\Gamma$ --- неспрямляемая кривая 


размерности Хаусдорфа $d_H>1$ (см., напр., \cite{18}, гл.\,3), то 


существуют нетривиальные 


функции, голоморфные в $\ov{\Bbb C}\setminus\Gamma$ 


и непрерывные в $\ov{\Bbb C}$. Известные 


примеры таких 


функций (в \cite{18}, \cite{19}) имеют вид 


\begin{equation} 


k(z)=\frac{1}{2\pi i}\int_E\frac{d\mu(\zeta)}{\zeta-z}, 


\end{equation} 


где $\mu$ --- мера Хаусдорфа на кривой $\Gamma$, 


а $E$ -- такое подмножество $\Gamma$, 


что $0<\mu(E)<\infty$. Ясно, что функция (17) является преобразованием 


Коши меры $\mu$, 


т.\,е. распределения, действующего по 


формуле 


\begin{equation} 


\mu(\omega)=\int_E\omega(\zeta)d\mu(\zeta). 


\end{equation} 


Интересно, что порядок сингулярности интегрирования $\mu$ с нулевой 


плотностью оказывается равным нулю, т.\,е. оно продолжимо на  


пространство всех непрерывных на $\Gamma$ функций, и в этом смысле 


ведет 


себя лучше ``естественного'' интегрирования (1). Обозначим через 


$W_\Gamma$ множество всех интегрирований с нулевой плотностью, заданных 


на $\Gamma$, а через $\wt W_\Gamma$ --- множество их преобразований Коши. 


Если $\Gamma$ --- 


замкнутая кривая, то $\wt W_\Gamma$ состоит из всех голоморфных в 


$\ov{\Bbb C}\setminus\Gamma$ и 


непрерывных в $\ov{\Bbb C}$ 


функций, обращающихся в нуль в точке $\infty$; если 


же $\Gamma$ --- разомкнутая дуга, то входящие в $\wt W_\Gamma$ 


функции могут иметь интегрируемые особенности на концах $\Gamma$. Таким 


образом, на замкнутой 


кривой множество $\wt W_\Gamma$ имеет структуру банаховой алгебры. В 


50-х\,--\,60-х годах Вермер рассматривал алгебру голоморфных в 


$\ov{\Bbb C}\setminus\Gamma$ и непрерывных в $\ov{\Bbb C}$ функций для 


случая, когда $\Gamma$ --- дуга ненулевой 


площади. Многими свойствами алгебры Вермера обладают и алгебры 


$\wt W_\Gamma$. 


Так, повторяя рассуждения из (\cite{20}, с.\,46--47), нетрудно 


показать, что если множество $\wt W_\Gamma$ не сводится к 


тождественному нулю, то оно содержит три функции, разделяющие точки 


$\ov{\Bbb C}$, и любая функция из этого множества принимает на $\Gamma$ 


все значения, которые она 


принимает в $\ov{\Bbb C}$. 


 


Если $\varphi$ есть интегрирование на $\Gamma$ с плотностью $f$ и 


$\psi\in W_\Gamma$, 


то $\varphi+\psi$ есть интегрирование с той же плотностью, так что при 


$W_\Gamma\ne\{0\}$ (в частности, при $d_H>1$) интегрирование 


восстанавливается по заданной 


плотности неединственным образом. Однако из множества всех  


интегрирований нередко можно выбрать одно, обладающее каким-либо  


дополнительным свойством. Так, согласно теореме Е.П.\,Долженко \cite{19} 


всякая функция $F\in H_\mu(\ov U)$ (где $U\supset\Gamma$ 


 --- область в $\Bbb C$), 


голоморфная в $U\setminus\Gamma$, 


должна быть голоморфной в $U$ при $\mu>d_H-1$. В связи 


с этим введем в рассмотрение следующие классы голоморфных функций. 


Пусть $0<\mu\le 1$. Будем относить голоморфную в 


$\ov{\Bbb C}\setminus\Gamma$ функцию $F$ к 


классу $\cal H_\mu(\Gamma)$, 


если у любой точки $t\in\Gamma$ (за исключением концов $\Gamma$, 


если $\Gamma$ --- разомкнутая дуга) имеется такая окрестность $U_t$ в 


$\Bbb C$, что $F$ удовлетворяет условию Г\"ельдера с показателем $\mu$ в 


замыкании 


любой связной компоненты $U_t\setminus\Gamma$. Пусть $\Gamma$ --- 


замкнутая кривая, тогда $F\in\cal H_\mu(\Gamma)$, если 


$F|_{D^+}\in H_\mu(\ov{D^+})$ и $F|_{D^-}\in H_\mu(\ov{D^-})$. 


Из теоремы Е.П.\,Долженко немедленно следует 


 


\begin{thm} Если $\mu>d_H-1$, то существует не более одного 


интегрирования с заданной плотностью, преобразование Коши которого 


принадлежит классу $\cal H_\mu(\Gamma)$. 


\end{thm} 


 


В частности, функция (11) принадлежит классу $\cal H_\mu(\Gamma)$ 


при $\mu<1-2p^{-1}_1(\nu,d)$. 


Это следует из классических оценок интеграла, 


входящего в формулу (11) (см., напр., \cite{21}). Отсюда получаем 


 


\begin{cor} 


Пусть выполнены условия теоремы 1 и, кроме того, 


\begin{equation} 


d_H-1<\mu<\frac{2\nu-d}{2-d}. 


\end{equation} 


Тогда существует единственное интегрирование с плотностью 


$f\in H_\nu(\Gamma)$, 


преобразование Коши которого принадлежит классу $\cal H_\mu(\Gamma)$. 


\end{cor} 


 


Кроме вопроса о единственности интегрирования с заданной плотностью, 


возникает вопрос о единственности интеграла, т.\,е. о выполнении 


равенства 


\begin{equation} 


\varphi_1(g_1)=\varphi_2(g_2) 


\end{equation} 


при условии, что $\varphi_j$ есть интегрирование с плотностью $f_j$, 


$j=1,2$, 


и 


\begin{equation} 


f_1(t)g_1(t)=f_2(t)g_2(t),\quad t\in\Gamma. 


\end{equation} 


Как мы только что видели, равенство (20) может быть нарушено даже 


при $f_1=f_2$, $g_1=g_2$, поскольку условие $f_1=f_2$ не влечет, 


вообще говоря, совпадения 


интегрирований $\varphi_1$ и $\varphi_2$. Однако сейчас мы покажем, что 


условие (21) 


все же влечет равенство (20), если $\varphi_{1,2}$ --- 


это те интегрирования, о которых идет речь в теореме 3 и следствии 2. 


Преобразования Коши этих интегрирований имеют вид (11). Подставив 


эти функции в (15), после стандартных преобразований получим 


(20). Таким образом, справедливо 


 


\begin{cor} 


Пусть $\Gamma$ есть кривая клеточной размерности $d$ и 


размерности Хаусдорфа $d_H$, а $f_j$, $g_j$ --- заданные на ней 


функции, $j=1,2$. Если $f_j\in H_\nu(\Gamma)$, 


$g_j\in H_\lambda(\Gamma)$, $j=1,2$, причем 


числа $\nu$ и $\lambda$ удовлетворяют 


условиям $\nu>d/2$, $\lambda>d-1$ и (19) (в случае, когда $\Gamma$ есть 


разомкнутая дуга, то сюда нужно добавить условия (12) и (16)\,), а 


$\varphi_j$ --- интегрирование с плотностью $f_j$ 


и преобразованием Коши из класса $\cal H_\mu(\Gamma)$, то из равенства 


(21) следует (20). 


\end{cor} 


 


В частности, полагая $f_2\equiv 1$, убеждаемся, что вычисление значений 


интегрирования, о котором идет речь в следствии 2, сводится к 


вычислению интегрирований с единичной плотностью (1), (4), (5). 


 


\section{Обобщенные меры} 


 


Наше определение интегрирования основано на формуле (9), позволяющей 


восстановить по интегрированию его плотность. Однако было 


бы интересно уметь восстанавливать плотность посредством предельного 


перехода типа дифференцирования интеграла по верхнему пределу. Для 


этого нужно определить интегрирования по частям $\gamma$ кривой 


$\Gamma$, которые могут затем стягиваться в точку. Если $\varphi$ --- 


интегрирование 


на $\Gamma$, $U\subset\Bbb C$ есть открытое множество и 


$\gamma=\Gamma\cap U$, то 


можно определить интегрирование по $\gamma$ как сужение $\varphi$ на $U$. 


Но тогда мы можем применять $\varphi$ только к функциям с носителями в 


$U$, т.\,е. 


нельзя применить последовательность таких сужений, соответствующих 


стягивающейся в точку последовательности множеств $U$, к одной и 


той же пробной функции $\omega$. Значит, желая осуществить 


вышеописанный 


предельный переход, мы должны с самого начала рассматривать не 


одно распределение $\varphi$, а семейство распределений, каждое из 


которых 


представляет собой интегрирование по некоторой дуге 


$\gamma\subset\Gamma$. 


 


\begin{defn} Система дуг $a(\Gamma)$ состоит из всевозможных дуг 


кривой $\Gamma$, направление обхода которых индуцировано ориентацией 


$\Gamma$. Она содержит как замкнутые (т.\,е. содержащие свои концы) дуги, 


так 


и открытые (т.\,е. не содержащие своих концов) и полуоткрытые 


(т.\,е. содержащие один из концов). С каждой из дуг 


$\gamma\in a(\Gamma)$ мы связываем ее 


начало $b(\gamma)$, конец $e(\gamma)$ и хорду 


$c(\gamma)=e(\gamma)-b(\gamma)$. 


\end{defn} 


 


\begin{defn} Обобщенной мерой на $\Gamma$ будем называть 


отображение $\varphi:a(\Gamma)\to\cal D'(\Bbb C)$, ставящее в 


соответствие каждой дуге $\gamma\in a(\Gamma)$ 


обобщенную функцию $\varphi(\gamma,\cdot)$ таким образом, что: 


\begin{itemize} 


\item[a)] $\supp\varphi(\gamma,\cdot)\subset\ov\gamma$; 


\item[б)] если $\gamma_1,\gamma_2,\dotsc$ --- конечная или счетная 


система попарно непересекающихся дуг, составляющих в совокупности одну 


дугу $\gamma$, то 


\begin{equation} 


\varphi(\gamma,\cdot)=\varphi(\gamma_1,\cdot)+ 


\varphi(\gamma_2,\cdot)+\dotsb; 


\end{equation} 


\end{itemize} 


если система $\{\gamma_j\}$ счетна, то равенство (22) означает, в 


частности, 


что ряд в его правой части сходится. 


\end{defn} 





С каждой точкой $t\in \Gamma$ свяжем однопараметрическое семейство дуг 


$\gamma_{t,\varepsilon}$. При фиксированном $\varepsilon>0$ дуга 


$\gamma_{t,\varepsilon}$ представляет 


собою связную компоненту пересечения 


$\Gamma\cap\{z:|t-z|\le\varepsilon\}$, содержащую точку $t$. 


 


\begin{defn} 


Будем говорить, что обобщенная мера $\varphi$ имеет 


плотность $f=f(t)$ в точке $t\in\Gamma$, если для любой функции 


$\omega\in C^\infty_0$ 


существует предел 


\begin{equation} 


\lim_{\varepsilon\to 0}\frac{1}{c(\gamma_{t,\varepsilon})} 


\varphi(\gamma_{t,\varepsilon},\omega)=f(t)\omega(t). 


\end{equation} 


Иначе говоря, обобщенная мера $\varphi$ 


имеет плотность $f(t)$, если 


\begin{equation} 


\lim_{\varepsilon\to 0}\frac{1}{c(\gamma_{t,\varepsilon})} 


\varphi(\gamma_{t,\varepsilon},\cdot)=f(t)\delta_t. 


\tag{23${}^1$}


\end{equation} 


\end{defn} 


 


Теперь приведем 


некоторые примеры обобщенных 


мер и вычислим их плотности. 


 


\section{Примеры} 


 


1. Простейшим примером обобщенной меры является обычная мера на 


$\Gamma$. Если $\mu$ есть заданная на $\Gamma$ конечная мера,то ее можно 


отождествить с обобщенной мерой $\varphi_\mu$, действующей по 


правилу 


\begin{equation} 


\varphi_\mu(\gamma,\omega)=\int_\gamma \omega\,d\mu. 


\end{equation} 


Если $\Gamma$ есть самоподобный фрактал размерности $d$, $1<d<2$, то 


$\Gamma$ есть $d$-множество (\cite{10}, \cite{22}), т.\,е. $d$-мерная 


мера Хаусдорфа $\mu$ удовлетворяет неравенству 


$$ 


C^{-1}r^d\le\mu(\Gamma\cap\{z:|z-t|\le r\})\le C r^d, 


$$ 


где постоянная $C$ не зависит ни от $t\in \Gamma$, ни от $r\in(0,1]$. 


В частности, $\mu(\gamma_{t,\varepsilon})\le C\varepsilon^d$. 


Для многих классических фракталов 


$c(\gamma_{t,\varepsilon})\ge c'\varepsilon$, где $c'$ не зависит от 


$\varepsilon$. Тогда 


предел (23) есть нуль, т.\,е. соответствующая обобщенная мера имеет 


плотность нуль в каждой точке кривой. Напомним, что порождаемое мерой 


Хаусдорфа интегрирование (18) имеет плотность нуль в смысле 


предыдущего параграфа. 


 


Очевидно, условие аддитивности (22) позволяет продолжить всякую 


обобщенную меру с семейства дуг $a(\Gamma)$ на множество 


$a_{\text{fin}}(\Gamma)$, содержащее всевозможные конечные объединения дуг 


семейства $a(\Gamma)$. Обсуждаемые в 


данном пункте обычные меры продолжимы, кроме того, на 


$\sigma$-алгебру $a_\sigma(\Gamma)\supset a_{\text{fin}}(\Gamma)$, 


содержащую все счетные объединения дуг. В 


следующем пункте рассмотрим меры, не обладающие свойством 


$\sigma$-продолжимости. 


\medskip 


 


2. С каждой дугой $\gamma\in a(\Gamma)$ свяжем логарифмическое ядро 


\begin{equation} 


k_\gamma(z)=\frac{1}{2\pi i}\ln\frac{z-b(\gamma)}{z-e(\gamma)}, 


\end{equation} 


где ветвь логарифма выделена посредством разреза по дуге $\gamma$ и 


условия 


$k_\gamma(\infty)=0$. Если эта функция $k_\gamma$ интегрируема в любой 


конечной части плоскости 


для любой дуги $\gamma$, то запишем 


обобщенную меру 


\begin{equation} 


\varphi(\gamma,\omega)=-\iint\limits_{\Bbb C} 


\myfrac{\omega}{\ov z}k_\gamma(z)dz\,d\ov z. 


\end{equation} 


Условие интегрируемости выполнено, в частности, если каждая дуга 


$\gamma\subset\Gamma$ является достижимой, т.\,е. если точки $b(\gamma)$ 


и $e(\gamma)$ можно соединить спрямляемой дугой $\gamma'$, не имеющей с 


$\gamma$ иных общих точек. При этом дуги 


$\gamma$ и $\gamma'$ в совокупности ограничивают область $D(\gamma)$. 


Если она расположена слева от $\gamma$, то обобщенную меру (26) 


можно записать в форме 


\begin{equation} 


\varphi(\gamma,\omega)=-\iint\limits_{D(\gamma)} 


\myfrac{\omega}{\ov z}dz\,d\ov z+\intl_{\gamma'}\omega\,dz;


\end{equation} 


в противном случае в этой формуле следует изменить знак. 


 


В терминах предыдущего параграфа можно сказать, что данная 


обобщенная мера связывает с каждой дугой $\gamma\subset\Gamma$ 


интегрирование 


по этой дуге с единичной плотностью. Покажем, что при некоторых  


геометрических ограничениях она имеет единичную плотность и в смысле 


определения 5. Действительно, 


$$ 


\int_{\gamma'}\omega\,dz=c(\gamma')w, 


$$ 


где $w$ лежит в замкнутой выпуклой оболочке множества 


$\omega(\gamma')$. 


Очевидно, $c(\gamma')=c(\gamma)$, 


а при стягивании дуги $\gamma$ в точку $t$ множество 


$\conv\omega(\gamma')$ 


стягивается в точку $\omega(t)$. Таким образом, 


$$ 


\lim_{\varepsilon\to 0}\frac{1}{c(\gamma_{t,\varepsilon})} 


\intl_{\gamma'_{\scriptstyle t,\varepsilon}}\omega(z)dz=\omega(t), 


$$ 


и нам осталось установить, что предел 


$$ 


\lim_{\varepsilon\to 0}\frac{1}{c(\gamma_{t,\varepsilon})} 


\iint\limits_{D(\gamma_{\scriptstyle t,\varepsilon})}


\myfrac{\omega}{\ov z} dz\,d\ov z 


$$ 


есть нуль. В качестве замыкающей дуги $\gamma'_{t,\varepsilon}$ можем 


взять дугу окружности $|z-t|=\varepsilon$; тогда 


$D(\gamma_{t,\varepsilon})$ есть часть соответствующего круга, и двойной 


интеграл имеет порядок $\varepsilon^2$. Таким образом, обобщенная 


мера (26) имеет плотность единица, если 


\begin{equation} 


\lim_{\varepsilon\to 0} 


\frac{\varepsilon^2}{c(\gamma_{t,\varepsilon})}=0. 


\end{equation} 


Это условие выполнено для многих классических фракталов, таких, как 


снежинка Коха. 


 


Как уже отмечалось, любая обобщенная мера продолжима по аддитивности на 


конечные объединения дуг семейства $a(\Gamma)$. Однако обобщенную 


меру (26), вообще говоря, нельзя продолжить на произвольное счетное 


объединение дуг кривой $\Gamma$. Рассмотрим в качестве примера  


бесконечно-звенную ломаную $\Gamma$, построенную следующим образом. 


Пусть $\{x_j\}$ 


есть монотонно убывающая последовательность положительных 


чисел, $x_1=1$, $\lim x_j=0$. Положим 


$z_{2j-1}=x_{2j-1}$, $z_{2j}=(1+i)x_{2j}$, $j=1,2,3,\dotsc$ 


Ломаная $\Gamma$ состоит из прямолинейных отрезков, последовательно 


соединяющих точки $z_1,z_2,z_3,\dotsc,z_j,\dotsc$ 


Пусть функция $\omega$ тождественно равна единице в окрестности 


$\Gamma$. Тогда для 


любой дуги $\gamma\subset\Gamma$ имеем 


$\varphi(\gamma,\omega)=e(\gamma)-b(\gamma)$. 


Рассмотрим последовательность звеньев $\{\gamma_j\}$ 


ломаной $\Gamma$, взятых через один, т.\,е.{} $\gamma_j$ 


начинается в точке $z_{2j-1}$ и кончается в точке $z_{2j}$. Тогда 


$\varphi(\gamma_j,\omega)=z_{2j}-z_{2j-1}=ix_{2j}+(x_{2j}-x_{2j-1})$. 


Очевидно, ряд $\sum(x_{2j}-x_{2j-1})$ сходится 


всегда, но сходимость ряда $\sum x_{2j}$ равносильна конечности длины 


ломаной $\Gamma$. 


Таким образом, если $\Gamma$ имеет бесконечную длину, то ряд 


$\sum\varphi(\gamma_j,\omega)$ 


расходится и обобщенную меру объединения звеньев $\cup\gamma_j$ 


определить не удается. С другой стороны, меру $\varphi$ все же можно 


определить на многих счетных семействах дуг неспрямляемой кривой общего 


вида. Так, если $\Gamma$ --- замкнутая неспрямляемая кривая, а $\Delta$ 


--- область 


со спрямляемой границей, то обобщенную меру (27) нетрудно определить 


на множестве $\Gamma\cap\Delta$, даже если оно состоит из счетного числа 


дуг. Действительно, в этом случае $D^+\cap\Delta$ состоит из счетного 


числа подобластей 


$\Delta$, причем их совокупная граница состоит из $\Gamma\cap\Delta$ и 


системы $\gamma'$ 


граничных дуг $\Delta$ конечной суммарной длины. При их подходящей 


ориентации имеем 


$$ 


\varphi(\Gamma\cap\Delta,\omega)=-\iint\limits_{D^+\cap\Delta} 


\myfrac{\omega}{\ov z}dz\,d\ov z+\intl_{\gamma'}\omega\,dz. 


$$ 


Далее, пусть на $\Gamma$ задана функция $f\in H_\nu(\Gamma)$, $\nu>d/2$. 


Если логарифмическое ядро (25) любой дуги $\gamma\in a(\Gamma)$ 


интегрируемо в любой степени $p\ge 1$ (скажем, если любая 


дуга достижима), то мы можем рассмотреть 


семейство голоморфных функций 


$$ 


F_\gamma(z)=f^w(z)k_\gamma(z)-\frac{1}{2\pi i} 


\iint\limits_{\Bbb C}\myfrac{f^w}{\ov\zeta} 


\frac{k_\gamma(\zeta)d\zeta\,d\ov\zeta}{\zeta-z}, 


$$ 


где $f^w$ --- продолжение Уитни функции $f$ 


с компактным носителем, и 


обобщенную меру 


\begin{equation} 


\varphi_f(\gamma,\omega)=-\iint\limits_{\Bbb C}F_\gamma(z) 


\myfrac{\omega}{\ov z}dz\,d\ov z. 


\end{equation} 


Это --- интегрирование по $\gamma$ 


с плотностью $f$ 


в смысле предыдущего параграфа. По-видимому, при некоторых 


геометрических ограничениях типа (28) обобщенная мера (29) имеет 


плотность $f(t)$ и в смысле определения 5. Отметим, что задача о 


восстановлении обобщенной меры по ее 


плотности не может иметь единственного решения в виду наличия  


нетривиальных обобщенных мер нулевой плотности (см. предыдущий пункт). 


\medskip 


 


3. Существуют и другие содержательные примеры обобщенных мер. 


Так, согласно формуле Грина (напр., \cite{23}, с.\,39) в области 


$D^+$ со спрямляемой границей $\Gamma$ справедливо соотношение 


$$ 


\intl_\Gamma u\myfrac{v}{n}ds=-\intl_{D^+} 


\biggl(\myfrac{u}{x}\myfrac{v}{x}+\myfrac{u}{y}\myfrac{v}{y} 


+u\nabla^2v\biggr)dx\,dy, 


$$ 


где $u$, $v$ --- заданные в $\ov{D^+}$ функции, $\nabla^2$ 


--- лапласиан, $ds$ --- элемент 


длины, $\myfrac{}{n}$ -- нормальная производная. Пусть теперь $\gamma$ 


--- дуга $\Gamma$, 


а $\gamma'$ --- лежащая в $D^+$ спрямляемая дуга с теми же началом и 


концом. 


Если $\gamma$ и $\gamma'$ в совокупности ограничивают область 


$D(\gamma)$, то правая 


часть формулы 


$$ 


\intl_\gamma u\myfrac{v}{n}ds=-\iint\limits_{D(\gamma)} 


\biggl(\myfrac{u}{x}\myfrac{v}{x}+\myfrac{u}{y}\myfrac{v}{y} 


+u\nabla^2v\biggr)dx\,dy+\intl_{\gamma'} 


u\myfrac{v}{n}ds 


$$ 


имеет 


смысл независимо от того, спрямляема дуга $\gamma$ или нет. Таким 


образом, на неспрямляемой кривой $\Gamma$ можем определить обобщенные 


меры 


\begin{align} 


\varphi_v(\gamma,\omega)&=-\iint\limits_{D(\gamma)} 


\biggl(\myfrac{\omega}{x}\myfrac{v}{x}+\myfrac{\omega}{y} 


\myfrac{v}{y}+\omega\nabla^2v\biggr)dx\,dy+ 


\intl_{\gamma'}\omega\myfrac{v}{n}ds,\\ 


% 


\psi_u(\gamma,\omega)&=-\iint\limits_{D(\gamma)} 


\biggl(\myfrac{u}{x}\myfrac{\omega}{x}+\myfrac{u}{y} 


\myfrac{\omega}{y}+u\nabla^2\omega\biggr)dx\,dy+ 


\intl_{\gamma'}u\myfrac{\omega}{n}ds. 


\end{align} 


В соотношениях (30), (31) 


функции $u(x,y)$ и $v(x,y)$ фиксированы, причем 


$v\in C^2(\ov{D^+})$, $u\in C^1(\ov{D^+})$. 


Обе эти обобщенные меры могут рассматриваться как усредненные значения 


нормальной производной (соответственно, фиксированной функции $v$ или 


пробной функции $\omega$) на дуге $\gamma$. 


При этом на фрактальной кривой $\Gamma\supset\gamma$ не существует в 


классическом смысле 


ни элемента длины, ни нормали. Таким образом, понятие 


обобщенной меры позволяет объединить в рамках единой концепции 


обычные меры, обобщенные интегралы на фрактальной кривой и другие 


содержательные объекты. Это дает основания полагать, что такие меры 


могут стать полезным инструментом фрактального анализа. 


 


 


\begin{thebibliography}{99} 


 


\bibitem{1} 


Mandelbrot B.B {\em The fractal geometry of nature}. -- San Francisco: 


Freeman, 1982.


\bibitem{2} 


Федер Е. {\em Фракталы}. -- М.: Мир, 1991. -- 260 с. 


\bibitem{3} 


Кац Б.А. {\em Задача о скачке и интеграл по неспрямляемой кривой} // 


Изв. вузов. Математика. -- 1987. -- \No\,5. -- С.\,49--57.


\bibitem{4} 


Harrison J., Norton A. {\em Geometric integration on fractal curves in 


the plane} // Indiana Univ. Math. J. -- 1991. -- V.\,40. -- \No\,2. -- 


P.\,567--594. 


\bibitem{5} 


Harrison J. {\em Stokes' theorem for nonsmooth chains} // Bull. Amer. 


Math. Soc. --  


1993. -- V.\,29. -- \No\,2. -- P.\,235--242. 


\bibitem{6} 


Стейн И. {\em Сингулярные интегралы и дифференциальные свойства 


функций}. -- М.: Мир, 1973. -- 342 с. 


\bibitem{7} 


Harrison J., Norton A. {\em The Gauss--Green theorem for fractal 


boundaries} // Duke Math. J. --  


1992. -- V.\,67. -- \No\,3. -- P.\,575--586. 


\bibitem{8} 


Кац Б.А. {\em Об интегрировании по неспрямляемой кривой} // Вопр. мат., 


мех. сплош. сред и применения мат. методов в стр-ве. -- М.: МИСИ,  


1992. -- С.\,63--69. 


\bibitem{9} 


Кац Б.А. {\em Об одной версии краевой задачи Римана на фрактальной 


кривой} // Изв. вузов. Математика. -- 1994. -- \No\,4. -- С.\,10--20. 


\bibitem{10} 


Jonssen A., Wallin H. {\em Function spaces on subsets of $\Bbb R$}. -- 


Math. 


Reports 2, Part 1, Harwood Acad. Publ. -- London, 1984.


\bibitem{11} 


Объедков Т.Е. {\em Применение нестандартного анализа к теории  


интегрирования и решение задачи Римана о скачке} // Тез. докл.


междунар. конф. ``Алгебра и анализ''. Ч.\,II. --


Казань, 1994. -- C.\,96--97


\bibitem{12} 


Хермандер Л. {\em Анализ линейных дифференциальных операторов с  


частными производными}. Т.\,1. -- М.: Мир, 1986. -- 462 с. 


\bibitem{13} 


Дынькин Е.М. {\em Гладкость интегралов типа Коши} // Зап. 


научн. семин. ЛОМИ АН СССР. -- 


1979. -- Т.\,92. -- С.\,115--133. 


\bibitem{14} 


Гахов Ф.Д. {\em Краевые задачи}. -- 3-е изд. 


-- М.: Наука, 1977. -- 640 с. 


\bibitem{15} 


Мусхелишвили Н.И. {\em Сингулярные интегральные уравнения}. -- 3-е 


изд., испр. и доп. -- М.: Наука, 1968. -- 511 с. 


\bibitem{16} 


Кац Б.А. {\em Задача Римана на замкнутой жордановой кривой} // Изв. 


вузов. Математика. -- 1983. -- \No\,4. -- С.\,68--80. 


\bibitem{17} 


Кац Б.А. {\em Задача Римана на разомкнутой жордановой кривой} // Изв. 


вузов. Математика. -- 1983. -- \No\,12. -- С.\,30--38. 


\bibitem{18} 


Маркушевич А.И. {\em Избранные главы теории аналитических функций}. -- 


М.: Наука, 1976. -- 191 с. 


\bibitem{19} 


Долженко Е.П. {\em 0 ``стирании'' особенностей аналитических функций} // 


УМН. -- 1963. -- Т.\,18. -- \No\,4. -- С.\,135--142. 


\bibitem{20} 


Гамелин Т.В. {\em Равномерные алгебры}. -- М.: Мир, 1973. -- 334 с. 


\bibitem{21} 


Векуа И.Н. {\em Обобщенные 


аналитические функции}. -- 2-е изд., перераб. -- М.: Наука, 1988. -- 


509 с. 


\bibitem{22} 


V\"ais\"al\"a J., Vuorinen M., Wallin H. {\em Thick sets and 


quasi-symmetric maps}. -- Preprint 2. -- Reports of Dept of Math., Univ.


of Helsinki, 1992. -- 25 p. 


\bibitem{23} 


Хейман У., Кеннеди П. {\em Субгармонические функции}. T.\,1.


-- М.: Мир, 1980. -- 304 с. 


 


\end{thebibliography} 


 


\vskip 2cm 


 


\post{\it Казанская государственная}{\it Поступила} 


\post{\it архитектурно-строительная}{11.03.1996} 


\post{\it академия}{} 


 


\label{end} 


\end{document} 


